\section{Conclusion}
\label{sec:Conclusion}

We have presented EDG and EDG++, a framework that uses electron density as privileged information to improve geometry-based molecular models through cross-modal distillation. By representing ED as multi-view RGB-D images and training an ED-aware teacher, the approach eliminates quantum chemical calculations at inference. EDG++ adds reliability-aware selective distillation---motivated by a gradient noise analysis (Proposition~\ref{prop:gradient_noise}) identifying sufficient conditions for noise reduction---to avoid negative transfer from unreliable teacher predictions. % [Polish] Original: Experiments show improvements on 34/36 QM9 task-model combinations and 55/60 rMD17 energy/force predictions across three architectures, with the improvement hierarchy aligning with the Hohenberg-Kohn theorem: thermodynamic properties benefit most, followed by electronic structure properties.
Experiments demonstrate improvements on 34/36 QM9 task-model combinations and on 55/60 rMD17 energy and force predictions across three architectures, with an improvement hierarchy that aligns with the Hohenberg-Kohn theorem: thermodynamic properties benefit most, followed by electronic structure properties. We further demonstrate generalization to the 166-atom mini-protein Chignolin, far outside the ED teacher's training distribution. There, naive EDG measurably degrades the baseline ($+\!4.0\%$ energy, $+\!2.8\%$ force MAE), confirming the negative-transfer hypothesis in the wild, while EDG++'s reliability gating recovers a net improvement ($-\!6.8\%$ energy, $-\!2.9\%$ force MAE). The frozen reliability estimator additionally serves as an inference-time OOD detector during MD rollouts---a dual-use capability that emerges for free from our training-time design.

\textbf{Limitations.} (1)~All results are reported with a single random seed (seed=42). While the sign test ($p < 10^{-8}$ for 34/36 QM9 wins) and bootstrap CI analysis (Appendix~\ref{sec:appendix_bootstrap}) provide evidence that the improvements are consistent, multi-seed evaluation would strengthen the conclusions. (2)~EDG++ results are reported per-task-best, which provides an optimistic upper bound; the default-threshold configuration captures the majority of gains (Table~\ref{tab:incremental}). % [Polish] Original: (3)~On rMD17, SchNet and ViSNet use the fixed-threshold EDG++ configuration ($\kappa\!=\!0$) only; per-molecule $\kappa$ tuning is performed exclusively on SphereNet due to the small (50-sample) validation set, which limits the reliability of per-task tuning on additional architectures.
(3)~On rMD17, SchNet and ViSNet use only the fixed-threshold EDG++ configuration ($\kappa\!=\!0$); per-molecule $\kappa$ tuning is performed exclusively on SphereNet, because the small (50-sample) validation set limits the reliability of per-task tuning on additional architectures. (4)~The reliability estimator is frozen during downstream tasks; joint fine-tuning may yield further improvements. (5)~The out-of-distribution case study (Sec.~\ref{sec:case_study_chignolin}) is performed on a single biomolecule (Chignolin, 166 atoms) using a single student architecture (ViSNet) and short MD trajectories (\SI{50}{\pico\second}); broader validation across diverse biomolecular systems, additional student architectures, and longer simulations is left to future work.

\textbf{Future Directions.} Several directions merit investigation: soft weighting schemes that modulate distillation strength continuously, molecule-adaptive distillation weights for conformationally challenging systems, and extension to other privileged information sources such as molecular orbitals or experimental spectra. The scale-invariant threshold $\tau = \mu + \kappa\sigma$ and the reliability-aware filtering mechanism can serve as general components for cross-modal knowledge transfer beyond molecular learning.

\textbf{Code Availability.} The source code and pre-trained models are available at \url{https://github.com/HongxinXiang/EDG}.
